% =========================================================
% 15. ensure_one() ORM Method
% =========================================================
\section{\texttt{ensure\_one()} ORM Method}

\begin{frame}[standout]
  \texttt{ensure\_one()} ORM Method
\end{frame}

\begin{frame}[fragile]{\texttt{ENSURE\_ONE()} METHOD}
  \begin{itemize}
    \item The \texttt{ensure\_one()} method verifies that a recordset contains
          \textbf{exactly one} record.
    \item If the recordset is empty or contains multiple records, it raises an error.
    \item It returns the same recordset when the check passes.
  \end{itemize}
  \begin{block}{Method Syntax}
\begin{lstlisting}
def ensure_one(self):
    return super().ensure_one()
\end{lstlisting}
    Important parts: \texttt{ensure\_one}, \texttt{self}.
  \end{block}
\end{frame}

\begin{frame}[fragile]{BREAKING DOWN THE METHOD SIGNATURE}
  \begin{itemize}
    \item \texttt{ensure\_one}: Safety check method for singleton recordsets.
    \item \texttt{self}: The recordset being validated.
    \item No extra arguments are required.
  \end{itemize}
\begin{lstlisting}
student = self.env['student.student'].browse(15)
student.ensure_one()
\end{lstlisting}
  \begin{itemize}
    \item This succeeds because there is exactly one record.
  \end{itemize}
\end{frame}

\begin{frame}[fragile]{WHAT DOES \texttt{ENSURE\_ONE()} RETURN?}
  \begin{itemize}
    \item On success, it returns the same recordset.
  \end{itemize}
\begin{lstlisting}
student = self.env['student.student'].browse(15)
result = student.ensure_one()
print(result)
\end{lstlisting}
  Output is:
\begin{lstlisting}
student.student(15,)
\end{lstlisting}
  \begin{itemize}
    \item On failure, it raises \texttt{ValueError}.
  \end{itemize}
\end{frame}

\begin{frame}[fragile]{WHEN IT FAILS}
  Empty recordset:
\begin{lstlisting}
students = self.env['student.student'].browse([])
students.ensure_one()  # ValueError
\end{lstlisting}
  Multiple records:
\begin{lstlisting}
students = self.env['student.student'].browse([15, 16])
students.ensure_one()  # ValueError
\end{lstlisting}
  \begin{itemize}
    \item Error message idea: ``Expected singleton''.
  \end{itemize}
\end{frame}

\begin{frame}[fragile]{WHY IS \texttt{ENSURE\_ONE()} IMPORTANT?}
  \begin{itemize}
    \item Many methods are written for one record only (send email, print report, open wizard).
    \item Accessing \texttt{record.name} on a multi-recordset can raise singleton errors later.
    \item Calling \texttt{ensure\_one()} makes the expectation explicit and fails early.
  \end{itemize}
\begin{lstlisting}
def action_confirm(self):
    self.ensure_one()
    self.state = 'confirmed'
    return True
\end{lstlisting}
\end{frame}

\begin{frame}[fragile]{METHOD CALL}
  \textbf{Method Call}
\begin{lstlisting}
self.ensure_one()
\end{lstlisting}
  \begin{itemize}
    \item Here, \texttt{.ensure\_one()} is the method call.
    \item Use it at the start of button actions and record-specific business methods.
  \end{itemize}
\end{frame}
